\documentclass[UTF8,a4paper,12pt]{ctexart}

\usepackage[a4paper,left=2.6cm,right=2.6cm,top=2.5cm,bottom=2.5cm]{geometry}
\usepackage{booktabs}
\usepackage{array}
\usepackage{enumitem}
\usepackage{xcolor}
\usepackage{hyperref}
\usepackage{fancyhdr}
\usepackage{microtype}
\usepackage{longtable}
\usepackage{tabularx}
\usepackage{listings}

\definecolor{LearnTraceBlue}{RGB}{30,83,120}
\definecolor{LearnTraceGray}{RGB}{90,98,108}
\definecolor{LearnTraceLight}{RGB}{238,245,247}
\hypersetup{
  colorlinks=true,
  linkcolor=LearnTraceBlue,
  urlcolor=LearnTraceBlue,
  pdfauthor={LearnTrace 项目组},
  pdftitle={LearnTrace 作品简介}
}

\setlength{\parindent}{2em}
\setlength{\parskip}{0.35em}
\setlist[itemize]{leftmargin=2.2em,itemsep=0.25em,topsep=0.35em}
\setlist[enumerate]{leftmargin=2.4em,itemsep=0.25em,topsep=0.35em}
\renewcommand{\arraystretch}{1.35}
\setcounter{tocdepth}{2}
\lstset{
  basicstyle=\ttfamily\small,
  columns=fullflexible,
  frame=single,
  rulecolor=\color{LearnTraceGray!35},
  backgroundcolor=\color{LearnTraceLight},
  breaklines=true,
  showstringspaces=false,
  xleftmargin=0.5em,
  xrightmargin=0.5em
}

\pagestyle{fancy}
\fancyhf{}
\fancyhead[L]{\small LearnTrace 作品简介}
\fancyhead[R]{\small AI 协作学习档案智能体}
\fancyfoot[C]{\thepage}
\setlength{\headheight}{15pt}

\title{\vspace{-1.5cm}\textbf{LearnTrace：面向 AI 辅助课程项目的\\可追溯学习档案智能体}\\[0.6em]
\large 作品简介}
\author{LearnTrace 项目组}
\date{2026 年 9 月}

\begin{document}

\maketitle
\thispagestyle{fancy}

\begin{center}
\small
\begin{tabular}{rl}
\textbf{作品类别：} & 可部署智能体 / 本地学习档案工具 \\
\textbf{项目版本：} & LearnTrace 0.1.0 \\
\textbf{项目仓库：} & \url{https://github.com/constantinmay/learntrace-skill} \\
\textbf{文档状态：} & 参赛提交版（2026 年 9 月 6 日）
\end{tabular}
\end{center}

\begin{abstract}
生成式人工智能正在成为编程学习和课程项目中的日常工具。面对这一变化，我们认为教育评价不应继续执着于分辨学生是否使用了 AI，而应把注意力转向一个更有价值的问题：学生在使用 AI 的过程中是否进行了思考、判断、修正和验证，又从中学到了什么。现有项目材料通常只能呈现最终结果，真正的学习过程则分散在聊天记录、Git 提交、设计文档和测试日志中，难以复盘，也难以被教师看见。

LearnTrace 是一个面向 AI 辅助课程项目的本地学习档案智能体。系统以\textbf{证据优先、学生确认、隐私最小化}为原则，在获得用户授权后，从本地 Git 仓库、任务书与设计文档、既有测试日志以及 OpenCode、Claude Code、Codex 等智能编码工具的会话导出中提取可追溯事件，构建候选学习节点，并通过结构化追问请学生确认、补充或否认，最终生成可编辑的学习档案和机器审计记录。作品不进行作弊检测、AI 作者归因、贡献排名或自动评分，而是帮助学生说明 AI 如何参与项目、自己如何作出决定，以及这些决定怎样经过实践验证。
\end{abstract}

\noindent\textbf{关键词：}AI 辅助学习；智能体；学习档案；证据追溯；Git；隐私保护；人机协作

\clearpage
\tableofcontents
\clearpage

\section{作品概览}

\begin{center}
\small
\begin{tabularx}{\textwidth}{>{\bfseries\color{LearnTraceBlue}}p{0.16\textwidth}X}
\toprule
核心问题 & AI 辅助开发留下大量分散痕迹，但最终作业通常只呈现代码和报告，学生的判断、修正与验证过程难以说明，教师也缺少可信的过程材料。 \\
目标用户 & 使用 AI 完成课程项目的学生，以及希望了解真实学习过程的教师和指导者。 \\
输入材料 & 本地 Git 历史、任务书与设计文档、已有测试日志，以及用户逐文件授权的 OpenCode、Claude Code、Codex 会话导出。 \\
核心处理 & 发现并授权范围，解析异构材料，建立证据索引，生成候选学习节点，向学生结构化追问，经确认后生成档案。 \\
主要输出 & 面向阅读的 \texttt{learning-record.md}、待确认问题清单和带来源索引与内容指纹的 JSON 审计档案。 \\
使用方式 & 可作为宿主 Agent Skill 调用本地 CLI，也可通过 \texttt{learntrace ui <project-dir>} 启动完整的本地 Web 智能体工作区。 \\
模型方式 & 核心证据工具链无需模型；交互式智能体支持 Anthropic Messages、OpenAI Chat Completions、OpenAI Responses 兼容 API，模型与网关由用户配置。 \\
完成状态 & 已形成可安装 Python wheel、内置 Web 前端和 Agent 服务；主分支持续集成覆盖 Ubuntu/Windows，测试、类型检查、代码规范和安装验证均已接入。 \\
\bottomrule
\end{tabularx}
\end{center}

\section{我们的理念}

AI 已经进入真实的学习过程。继续把“有没有使用 AI”作为首要问题，容易把学生和教师都带入一场很难可靠完成的检测：学生担心留下 AI 使用痕迹，教师只能从最终代码推测过程，而真正重要的学习行为反而被忽略。我们不希望 LearnTrace 成为另一个侦测工具。

我们更关心的是：学生能否提出合适的问题，能否理解和检查 AI 的回答，能否在建议不适用时作出修改或拒绝，能否从失败中调整方案，并用测试、文档或代码变化验证自己的判断。使用 AI 本身既不能证明学生已经学习，也不应自动成为否定学习的理由。评价应当回到这些可以说明、可以核验的过程上。

因此，LearnTrace 不试图回答“这段代码到底是谁写的”，而是帮助学生回答“我如何借助 AI 推进项目，我在哪些地方作出了自己的判断，我从这些选择中学到了什么”。系统先整理可观察证据，再提出候选学习节点，最后把确认权交还给学生本人。教师看到的不是一个由模型给出的结论，而是一份可以继续追问、可以返回原始材料核验的学习档案。

这一理念也决定了系统的边界：

\begin{itemize}
  \item 不把 AI 使用视为需要隐藏的行为，而是鼓励透明、负责任地记录协作过程；
  \item 不依据 Git 或聊天记录推断作者身份，不输出贡献排名和学习分数；
  \item 不把系统推测包装成事实，证据不足时明确标注“未记录”或提出待确认问题；
  \item 不自动执行目标仓库的代码和测试，只读取用户选定的现有材料与日志；
  \item 不默认上传项目源码、完整聊天或个人路径，分析产物优先保存在本机。
\end{itemize}

我们希望由此建立一种更坦诚的人机协作方式：学生不必隐藏 AI 的参与，也不能把 AI 的输出直接当作自己的理解；教师无需执着于不可靠的 AI 检测，而可以关注学生有没有判断、修正、验证和反思。LearnTrace 记录的正是这一部分。

\section{作品总体设计}

LearnTrace 采用“确定性证据工具链 + 宿主智能体交互”的双层架构。Python 工具链负责文件发现、格式解析、证据索引、隐私脱敏、规则推断、Schema 校验和档案生成；宿主智能体负责理解用户意图、编排工具、展示候选学习节点并与学生完成确认。两层职责清晰分离：确定性的工作尽量由可测试的程序完成，需要语义理解的部分交由智能体完成，同时所有关键学习结论仍需用户确认。

\begin{center}
\small
\begin{tabular}{>{\raggedright\arraybackslash}p{0.20\textwidth}
                >{\raggedright\arraybackslash}p{0.32\textwidth}
                >{\raggedright\arraybackslash}p{0.36\textwidth}}
\toprule
\textbf{层次} & \textbf{主要模块} & \textbf{职责} \\
\midrule
输入与授权层 & 项目发现、范围选择、轨迹授权 & 先列出候选材料，再由用户确认读取范围；AI 会话导出逐文件授权 \\
证据解析层 & Git、文档、测试日志、轨迹适配器 & 把异构材料转换为带来源引用的统一可观察事件 \\
证据治理层 & 脱敏、校验、聚合、工作分段 & 删除敏感正文，限制数据规模，保证记录满足版本化 Schema \\
学习建模层 & 候选节点生成、学生确认 & 从事实中提出可质疑的候选关系，由学生确认、补充或否认 \\
产品交互层 & CLI、Skill、本地 Web UI、内置 Agent & 提供命令行和图形界面两种使用方式，支持流式对话与报告查看 \\
输出层 & Markdown 档案、问题清单、JSON 审计档案 & 同时满足人类阅读、后续编辑与机器审计需求 \\
\bottomrule
\end{tabular}
\end{center}

系统的核心数据流如下：

\begin{center}
\small
\texttt{授权输入}
$\Rightarrow$
\texttt{统一事实}
$\Rightarrow$
\texttt{候选学习节点}
$\Rightarrow$
\texttt{学生确认}
$\Rightarrow$
\texttt{可追溯档案}
\end{center}

在整个过程中，每一项关键内容都保留来源引用。系统可以指出结论来自哪一份文档、哪一条测试日志、哪一个 Git 对象或哪一段经授权的工具调用记录，使最终报告不只是自然语言总结，更是一份可以回到原始材料核验的证据地图。

\section{核心功能}

\subsection{项目材料发现与授权确认}

系统首先以只列目录、不读取正文的方式发现候选 Git 仓库、Markdown 或纯文本文档、测试日志和可选的 AI 轨迹文件，并把候选范围展示给用户。只有在用户确认后，系统才进入实际解析阶段。对于可能包含完整聊天、工具参数和源码片段的 AI 会话导出，LearnTrace 要求逐个文件显式授权；一个文件的授权不会自动扩展到其他会话。这种“先发现、后读取”的流程兼顾了易用性与知情同意。

\subsection{多源项目证据解析}

LearnTrace 能够读取本地 Git 历史、任务书与设计报告、已有测试输出及其他运行日志。Git 模块不仅提供提交摘要，还能够生成完整但轻量的历史索引，记录提交树、文件角色、源码与测试是否共同出现等事实；当智能体需要理解某次关键修改时，再按提交、路径、行范围或字节范围回读证据，避免一次性加载大型仓库。系统也支持检查暂存区和当前工作区，使尚未提交但确实存在的阶段性成果能够被如实记录。

对于二进制文件、Git LFS 指针、子模块、非 UTF-8 文本或超长内容，系统不会强行解释，而是保留类型、对象标识、大小和不可读原因。对较长 Git 历史和会话文件采用索引、分页、流式解析与数量上限，在保持可追溯性的同时控制内存和上下文开销。

\subsection{多宿主 AI 协作轨迹适配}

作品实现了 OpenCode、Claude Code 和 Codex 三类会话导出的适配器，将不同格式的工具调用统一转换为版本化的可观察事件。系统主要保留时间、工具类型、项目内相对路径、保守的命令类别、执行状态和来源宿主；不把完整提示词、隐藏推理、聊天正文、工具输出、补丁正文、密钥和个人绝对路径写入公开档案。

适配器只在观察到工具返回结果时记录完成事件；非零退出码会被标为错误，状态无法确认时会标为未知，未找到对应返回的调用则被跳过并产生警告。对于跨越多个会话的项目，系统可合并事件、稳定去重并保留原会话标识，还能依据时间间隔和明确的项目对象范围进行确定性工作分段，从而形成连续但不过度推断的协作过程。

\subsection{证据分层与候选学习节点}

LearnTrace 将档案信息严格划分为三类：

\begin{enumerate}
  \item \textbf{可观察事实：}能够直接定位到 Git、文档、测试日志或授权轨迹；
  \item \textbf{学生确认：}学生针对候选节点给出的确认、补充或否认；
  \item \textbf{系统候选推断：}系统根据多个事实提出的可能解释，必须附带依据和不确定性。
\end{enumerate}

默认确定性推断器能够根据明确证据提出失败后修复、约束调整、补充测试和连续追问等候选，但不会仅凭时间接近或词语相似断言学生修改或拒绝了 AI 建议。后一类关系只有在宿主智能体获得完整会话读取授权并进行语义核对，或由学生本人补充确认时，才可进入候选或学生陈述。候选节点不是最终结论；系统会生成少量结构化问题，请学生说明当时的判断和收获。

\subsection{学习档案与审计产物}

一次完整分析会形成三类互补产物：

\begin{itemize}
  \item \texttt{learning-record.md}：供学生和教师阅读的精简学习档案，展示项目目标、阶段进展、AI 使用场景、关键决策、验证证据、个人反思和后续计划；
  \item \texttt{.learntrace/learning-questions.md}：集中保存仍需学生确认、补充或否认的问题；
  \item \texttt{.learntrace/archive-records.json}：保存完整事实、来源索引、候选关系、警告和内容指纹，支持机器审计与结果复现。
\end{itemize}

确认阶段复用首次分析生成的事实快照，不重新读取轨迹或偷偷改变候选依据。Web 界面还为每次会话保存不可变的阶段产物快照，因此后续重新分析同一项目时，旧会话中的报告和证据引用不会被新结果篡改。

\subsection{本地 Web 智能体工作区}

除可安装的 \texttt{learntrace} CLI 与 Agent Skill 外，作品提供本地 Web 工作区。用户执行 \texttt{learntrace ui <project-dir>} 后即可在浏览器中完成模型设置、项目选择、普通对话、授权确认、结构化问答、进度查看、停止分析和报告阅读。界面通过服务器发送事件（SSE）流式呈现消息、工具活动和产物状态；会话元数据保存在本地 SQLite 中，服务重启后仍可恢复同一智能体上下文。

本地服务只监听回环地址，并通过一次性启动令牌和 HttpOnly Cookie 建立会话；模型 API Key 使用操作系统凭据库保存，不进入目标项目、浏览器存储、SQLite 数据库或普通 JSON 配置文件。同一项目同一时间只允许一个分析任务写入产物，降低并发覆盖风险。内置 Agent 的 \texttt{run\_command} 工具会校验命令和路径：核对过的 LearnTrace 操作可自动放行，授权、导出、范围外或其他命令必须请求用户批准。Agent 同时启用了只读的 \texttt{read/grep/find/ls} 工具，它们不经过该命令守卫；其使用范围由 Skill 流程和底层工具行为共同约束，因此本系统不把内置 Agent 宣称为操作系统沙箱。

\section{模型与 API 调用方式}

LearnTrace 不绑定某一家模型，也不把模型能力混入所有底层步骤，而是采用\textbf{单智能体、可配置模型、本地确定性工具}的组合方式。

\begin{itemize}
  \item \textbf{CLI 路径：}证据解析、脱敏、候选生成、Schema 校验和报告渲染全部由本地 Python 程序以确定性规则完成，不调用远程大模型，也不要求 API Key。这使核心结果可测试、可重复，并可在离线环境中运行。
  \item \textbf{Skill 路径：}当 LearnTrace 作为 Codex、Claude Code、OpenCode 等宿主中的 Skill 使用时，模型登录、API 连接和对话界面由宿主负责；LearnTrace CLI 只承担本地证据工具的角色。
  \item \textbf{内置 Web Agent 路径：}Web 工作区通过随包分发的 Pi SDK Agent 服务调用用户选择的模型。用户可配置 API 地址、模型 ID、API Key、协议和思考强度。当前支持 Anthropic Messages、OpenAI Chat Completions 和 OpenAI Responses 三类协议，因此既可接入相应云端服务，也可连接提供兼容接口的模型网关；本机回环地址还允许使用符合协议的本地服务。
\end{itemize}

模型在系统中的作用主要是理解用户自然语言、按照 Skill 编排工具、解释已结构化的证据以及发起必要追问。Git 读取、文件解析、隐私过滤、结果校验和档案落盘仍由确定性代码负责。API 调用只发生在用户主动配置并启动交互式智能体时，API Key 由操作系统凭据库管理；LearnTrace 的目标项目产物不保存密钥。这样的架构减少了多模型串联造成的信息漂移、成本叠加和责任不清，也使用户能够根据预算、隐私要求与可用服务自由选择模型。

\subsection{参赛演示配置与请求链路}

参赛演示记录使用学校提供的兼容模型网关，非敏感配置如下。模型名称是网关接受的运行时标识，不等同于对服务端底层实现的证明，也不会被硬编码进程序。

\begin{center}
\small
\begin{tabularx}{0.94\textwidth}{>{\bfseries}p{0.23\textwidth}X}
\toprule
项目 & 记录值 \\
\midrule
服务 & 中国科学技术大学兼容模型网关 \\
API 基址 & \url{https://api.llm.ustc.edu.cn/v1} \\
模型标识 & \texttt{qwen3.8-reasoner} \\
调用协议 & Anthropic Messages（\texttt{/v1/messages}） \\
验证口径 & 配置用于开发联调；现场可用性取决于网关状态、账号权限和当时返回结果 \\
\bottomrule
\end{tabularx}
\end{center}

更换其他可用模型时，只需在网页设置中修改 API 地址、协议和模型 ID，无需修改证据处理代码。API Key 仅由用户在本机界面填写并保存至操作系统凭据库，不写入源码、Git 历史、项目档案或本文档。

一次 Web 交互请求的实际链路为：

\begin{center}
\small
\texttt{浏览器（本机）}
$\rightarrow$
\texttt{FastAPI 会话服务}
$\rightarrow$
\texttt{内置 Pi SDK Agent}
$\rightarrow$
\texttt{用户配置的兼容模型 API}
\end{center}

模型返回工具调用意图后，内置 Agent 可以使用只读文件工具以及受控的 LearnTrace 命令工具；其中命令执行、授权和产物写入受到本地程序检查。选择远程 API 时，对话内容、模型所需的证据片段和工具结果可能进入远程模型上下文；“文件保存在本地”和“内容从不发送到远程模型”不是同一件事。若不希望发送语义内容到远程服务，用户应使用纯 CLI 路径或连接符合协议的本地模型服务。

\section{关键技术与实现难点}

\subsection{异构证据的统一语义}

Git 记录文件树和差异，测试日志记录执行结果，文档表达项目目标，智能编码工具又使用不同的 JSON 或 JSONL 会话格式。项目以版本化 Schema 将这些材料统一为带稳定标识、事件类型、摘要、时间和来源引用的记录。各适配器只负责格式转换，不直接生成学习结论，从架构上避免解析层越权推断。

\subsection{可追溯性与上下文规模的平衡}

完整复制源码、聊天和工具输出会造成上下文膨胀和隐私风险。LearnTrace 采用“轻量索引 + 按需回读”：先生成完整可检索索引，仅在核验关键修改时读取有限行数或字节范围。长会话逐行流式解析，并对文件、单行、事件和警告数量设置上限；截断会形成结构化警告，不会被悄悄隐藏。

\subsection{叙事质量与人机责任边界}

仅靠模型自动撰写“学习收获”容易产生事后合理化。LearnTrace 先区分事实、候选和学生陈述，再通过结构校验、引用存在性检查和确认状态约束降低无依据叙述风险。ContractValidator 能检查可判定的契约错误，但不能证明每句自然语言都被引用证据完整支持；最终语义仍需学生或读者回到原始材料复核。

\subsection{本地智能体的安全运行}

本地服务使用回环监听、一次性令牌和 HttpOnly Cookie 保护接口；模型配置、对话元数据、项目产物和 API Key 分开存储。\texttt{run\_command} 命令守卫校验子命令、参数和项目写入范围，并提供超时、输出截断、取消及子进程清理。内置 \texttt{read/grep/find/ls} 是独立的只读工具，不经过同一命令守卫，因此系统明确将自身定位为有边界的本地工具，而不是完整操作系统沙箱。

\subsection{可恢复交互与产物一致性}

学习档案生成包含授权、工具调用、等待回答、继续分析和产物生成。作品通过 SQLite 会话状态、SSE 事件流、持久 Agent 会话、结构化问题卡片、项目写锁和不可变产物快照保持一致；刷新或重启后可以继续原会话，停止分析时则取消未完成问题并终止正在运行的任务。

\section{作品创新点}

\begin{enumerate}
  \item \textbf{证据分层与学生共同确认。} 系统把事实、候选推断和学生陈述分开保存，把学生设计为档案共同作者，而不是被模型分析和评分的对象。
  \item \textbf{隐私最小化的多来源过程重建。} Git、文档、测试日志和三类编码 Agent 轨迹被统一为可检索事件；确定性工具与模型解耦，既能离线复现，也避免仅凭时间邻近制造因果关系。
  \item \textbf{面向长周期项目的可回查闭环。} 从授权、索引、按需取证、追问到报告生成，关键内容保留来源；稳定去重、内容指纹和不可变快照使阅读档案与机器审计记录可以相互核对。
\end{enumerate}

\section{典型使用流程}

以一个使用 AI 完成课程编程项目的学生为例，完整流程如下：

\begin{enumerate}
  \item 学生在项目目录启动 LearnTrace CLI、宿主 Skill 或本地 Web 工作区；
  \item 系统发现 Git、任务书、设计报告和测试日志等候选材料，并请学生确认分析范围；
  \item 若学生希望纳入 AI 协作过程，则主动提供 OpenCode、Claude Code 或 Codex 会话导出，并逐份授权；
  \item 系统解析多源材料，生成可观察事件、Git 历史索引、轨迹工作分段和候选学习节点；
  \item 智能体围绕证据不足或需要本人判断的部分提出少量问题，学生进行确认、补充或否认；
  \item 系统生成面向阅读的学习档案、待确认问题清单和机器审计档案；
  \item 学生检查脱敏后的学习档案，将其用于个人复盘、课程总结或答辩材料，并妥善保留仅供本地审计的机器档案。
\end{enumerate}

该流程也支持“无 AI 轨迹”的降级模式：当学生没有提供或没有授权会话导出时，系统只根据 Git、文档与测试日志生成普通项目复盘，并明确标记 AI 协作信息未记录，不会反向猜测学生是否使用过 AI。

\section{脱敏案例：从失败日志到学习节点}

仓库中的 \texttt{examples/task4-demo} 提供了一组可独立检查的合成材料。案例模拟名单导入功能：设计文档要求支持带引号的 CSV 字段；已有 pytest 日志显示姓名“Smith, Jr.”被错误拆成两列；随后补丁把单一正则表达式替换为小型状态机，并增加引号逗号和空字段测试。

\begin{center}
\small
\begin{tabularx}{\textwidth}{>{\bfseries}p{0.17\textwidth}X}
\toprule
环节 & 案例中的实际内容 \\
\midrule
原始材料 & 设计约束、失败日志、脱敏会话片段和修复补丁均随示例提供 \\
系统事实 & 记录失败用例和后续提交，并保留到对应文件的来源引用 \\
候选问题 & “这次修改是否是为了修复日志里的失败？” \\
学生回答 & 示例确认文件明确标注为模拟回答，不冒充真实学生陈述 \\
最终档案 & 确认内容进入学习档案；未经授权容器绑定的 AI 建议记录被忽略并产生警告 \\
\bottomrule
\end{tabularx}
\end{center}

该案例证明的是证据、候选、确认和档案的工程闭环，不用于证明真实教学效果。默认推断器也不会把会话中的 AI 建议与后续提交自动配对；这项关系必须依赖完整授权语义或学生确认。

\section{技术栈与工程质量}

作品后端以 Python 3.11、FastAPI、Uvicorn、JSON Schema、SQLite、Keyring 等技术实现，负责 CLI、HTTP/SSE 服务、会话存储、证据解析和档案管理；前端采用 React、TypeScript 与 Vite 构建；内置智能体服务基于 Pi SDK，通过独立 Node.js 运行时与 Python 后端通信。项目提供可安装的 Python wheel，并把前端静态资源、Agent 服务和 Skill 一并打包，普通用户无需分别启动多个开发服务器。

工程侧建立了单元测试、集成测试、UI 测试和打包安装测试，覆盖静态解析、Git 导航、三类轨迹适配、隐私脱敏、Schema 校验、报告渲染、会话恢复、模型配置、命令守卫、进程取消与前端历史状态渲染等关键路径；同时使用 Ruff、Pyright、Pytest、TypeScript 类型检查和前端测试工具进行持续质量控制。系统对失败状态采取显式处理：缺失文件、未授权输入、格式错误、截断、无法识别的调用类型和未知执行状态都会形成可见结果或结构化警告。

最近一次公开持续集成基线为 Git 提交 \texttt{0af1135}。2026 年 9 月 6 日的 GitHub Actions 运行记录显示，Python 全量测试在 Ubuntu 和 Windows 上均为 \textbf{690 项通过、1 项跳过}，Web 测试 6 项通过，Agent 测试 81 项通过；Ruff、Pyright、前后端构建、资源同步和 wheel 检查同时通过。运行记录：\url{https://github.com/constantinmay/learntrace-skill/actions/runs/34011233296}。这些结果证明对应提交满足既定工程检查，但不替代对教学价值的案例观察和用户评价。

\begin{center}
\small
\begin{tabularx}{\textwidth}{>{\bfseries}p{0.19\textwidth}X X}
\toprule
验收层次 & 重点验证内容 & 代表性工具 \\
\midrule
代码与类型 & Python/TypeScript 格式、类型和静态缺陷 & Ruff、Pyright、TypeScript、ESLint \\
功能行为 & 解析器、适配器、Schema、报告、Web 会话与 Agent 协议 & Pytest、Vitest \\
集成链路 & 三类会话导入、归档闭环、授权边界和真实 Git 行为 & Python 集成测试 \\
发布交付 & 前端/Agent 资源同步、wheel 构建与隔离安装 & Vite、Rolldown、uv、pip \\
跨平台 & Windows 与 Linux 的路径、进程和打包行为 & GitHub Actions 双平台任务 \\
\bottomrule
\end{tabularx}
\end{center}

\section{安装部署与快速体验}

\subsection{运行条件}

核心 CLI 需要 Python 3.11 或更高版本；仅使用宿主 Skill 时无需另装 Node.js，使用自带 Web 智能体工作区则需要 Node.js 22.22.2 或更高版本。前端静态资源、Agent 服务和 Skill 已随 Python 包分发，普通用户无需进入源码目录重新安装前端依赖。

\subsection{最短使用路径}

\begin{lstlisting}
uv tool install .
learntrace --version
learntrace ui <project-dir>
\end{lstlisting}

页面打开后，用户依次完成模型设置、项目选择、材料范围确认、可选轨迹授权和结构化问题回答，即可查看学习档案。若只使用本地确定性管线，可以运行：

\begin{lstlisting}
learntrace discover <project-dir>
learntrace run <project-dir>
\end{lstlisting}

作品也可直接安装到 Claude Code：在安装 CLI 后，将完整的 \texttt{skills/learntrace/} 目录复制到项目内的 \texttt{.claude/skills/learntrace/}，或个人目录 \texttt{\textasciitilde/.claude/skills/learntrace/}，随后在待分析项目中输入：

\begin{lstlisting}
/learntrace 请分析当前项目并生成可追溯学习档案。
\end{lstlisting}

此时模型、对话与工具审批由 Claude Code 提供，LearnTrace 只调用本地确定性 CLI，不要求另配模型 API，也不需要 Node.js。直接调用 Skill 不会读取 Claude Code 历史；只有学生希望把 AI 协作过程作为证据时，才需自行复制指定会话 JSONL，并对该文件显式授权。未授权会话不影响 Git、项目文档和已有测试日志的正常分析。

\section{数据安全、适用边界与已知限制}

添加 AI 会话轨迹时必须逐份授权；LearnTrace 不主动扫描 Claude Code、Codex 或 OpenCode 历史目录，也不自动运行目标项目测试。机器审计档案可能包含本地来源索引，应被视为敏感产物；对外分享前应检查经过脱敏的主档案，不应直接上传原始会话或未经检查的机器归档。

当前版本主要支持 Git、Markdown/纯文本、常见测试日志和三类编码 Agent 会话格式；对 PDF、Word 等复杂二进制文档主要记录文件事实，不强行解释正文。“测试通过”必须来自用户已有日志，候选学习节点也不是自动评分，最终叙述仍需学生确认和复核。

这些限制是有意的产品边界。课程评价涉及真实个人经历，系统宁可明确缺失证据，也不为得到更完整的报告而扩大读取范围、自动执行不可信代码或让模型补写未经确认的经历。

\section{应用价值与后续发展}

LearnTrace 的直接价值是帮助学生把“做完项目”提升为“能够说明自己如何完成项目”。它能够减少期末整理材料的机械劳动，帮助学生定位真正值得写入报告和答辩的决策、失败与验证；同时，它为教师提供一种不同于代码查重或 AI 检测的观察视角，使教学关注点回到学生的判断、验证和反思。

后续版本可在现有证据模型上扩展更多文档格式和智能编码工具，增加面向团队项目的协作视图、课程模板与脱敏导出模板，并在获得明确授权的前提下支持更多学习管理场景。无论功能如何扩展，系统仍将坚持三个原则：不以推断替代事实，不以自动化替代学生主体性，不以便利为由突破隐私边界。

\section*{结语}

在 AI 已经成为学习工具的今天，简单追问“学生是否使用了 AI”越来越难以反映真实的学习情况。LearnTrace 希望把问题向前推进一步：不回避 AI 的参与，也不把使用 AI 等同于完成学习，而是认真记录学生怎样提问、怎样判断、怎样修正，以及怎样验证。

我们相信，真正值得被看见的不是一份看似完全由人独立完成的结果，而是学生在新工具参与之后仍然保有的思考和成长。LearnTrace 所做的，就是让这部分过程有证据可循、有本人确认，也有机会进入课程总结与评价。

\end{document}
